\section{Matroidal evidence closure}
\label{sec:matroid-evidence}

The four-valued kernel records whether positive and negative evidence are
present, but it does not by itself distinguish evidential novelty from coarse
repetition. This section adds a structural layer for that distinction. The
claim is deliberately restricted: arbitrary bodies of epistemic evidence are
\emph{not} assumed to form matroids. Instead, the bilateral 4-PEL signature
canonically induces a two-channel closure structure. The formal declarations
are in \texttt{PEL4/MatroidEvidence.lean}, with a separate declaration-level
trust inventory in \texttt{PEL4/MatroidEvidenceAxiomAudit.lean}.

\begin{definition}[Two-channel evidence closure]
Let \(E\) be a type of evidence atoms and let
\(\chi:E\to\{+,-\}\) assign each atom its coarse 4-PEL polarity. For
\(A\subseteq E\), define
\[
  \operatorname{cl}_{\chi}(A)
  = \{e\in E : \exists a\in A\;\chi(a)=\chi(e)\}.
\]
Thus, after coarse projection, one item of positive evidence generates the
whole positive parallel class, and likewise for negative evidence. No claim is
made that the underlying fine-grained sources are interchangeable.
\end{definition}

\begin{theorem}[4-PEL channel closure]
The operator \(\operatorname{cl}_{\chi}\) is extensive, monotone, idempotent,
and satisfies the closure exchange law
\[
 x\in \operatorname{cl}_{\chi}(A\cup\{y\}),\quad
 x\notin \operatorname{cl}_{\chi}(A)
 \quad\Longrightarrow\quad
 y\in \operatorname{cl}_{\chi}(A\cup\{x\}).
\]
Hence the coarse positive/negative partition carries a matroid closure
presentation. In particular,
\[
 y\in \operatorname{cl}_{\chi}(\{x\})
 \quad\Longleftrightarrow\quad
 \chi(x)=\chi(y).
\]
\textup{\textsc{Status: Lean-proved; separate declaration-level audit in
\texttt{PEL4/MatroidEvidenceAxiomAudit.lean}.}}
\end{theorem}
\begin{proof}[Proof sketch]
Extensivity and monotonicity follow immediately from retaining a witness of the
same polarity. Idempotence collapses a two-step same-polarity witness to a
one-step witness. For exchange, if \(x\) enters the closure only after inserting
\(y\), then the witnessing atom cannot already belong to \(A\); consequently
\(x\) and \(y\) have the same polarity, so inserting \(x\) places \(y\) in the
closure. This is the two-class partition instance of the standard closure
view of matroids \cite{Whitney1935,Oxley2011}.
\end{proof}

For a body of evidence \(A\), write
\(\operatorname{Supp}_{+}(A)\) and \(\operatorname{Supp}_{-}(A)\) for the
existence of positive and negative atoms respectively. These two predicates
recover exactly the bilateral Boolean observations used by the FDE kernel.

\begin{proposition}[Coarse-state conservativity]
For each polarity \(\sigma\in\{+,-\}\),
\[
  \operatorname{Supp}_{\sigma}(\operatorname{cl}_{\chi}(A))
  \quad\Longleftrightarrow\quad
  \operatorname{Supp}_{\sigma}(A).
\]
Consequently, if \(A\) realizes an FDE value by matching its positive and
negative bits to these two support predicates, then
\(\operatorname{cl}_{\chi}(A)\) realizes exactly the same value, and conversely.
Matroidal closure removes coarse redundancy without changing the four-valued
observation.
\textup{\textsc{Status: Lean-proved; separate declaration-level audit in
\texttt{PEL4/MatroidEvidenceAxiomAudit.lean}.}}
\end{proposition}

\begin{proposition}[Channel-rank profile]
Define the coarse channel rank of \(v\) as the number of present bilateral
support bits. Then
\[
  r(\Nval)=0,\qquad
  r(\Tval)=r(\Fval)=1,\qquad
  r(\Bval)=2.
\]
\textup{\textsc{Status: Lean-proved by reduction.}}
\end{proposition}

The rank profile changes the interpretation of contradiction in a useful way.
At this coarse structural level, \(\Bval\) is not deficient in evidence: it has
maximal channel rank because both independent directions are represented.
Paraconsistency is therefore compatible with high evidential breadth. Rank
alone is not a truth value, however. Since \(\Tval\) and \(\Fval\) both have
rank one, polarity orientation remains essential. The rank statistic measures
how many coarse evidence directions are represented, not which direction they
support and not how reliable or numerous their sources are.

\paragraph{Scope boundary.}
Same-polarity atoms are parallel only after projection to the two bits observed
by the FDE kernel. Two independent laboratories, two proofs, or two sensors may
remain genuinely distinct at a finer epistemic level even when both contribute
positive support. The present theorem therefore establishes a canonical
\emph{coarse} matroid induced by 4-PEL; it does not establish the general thesis
that evidence, justification, or information always has matroid structure.
Reliability tags, source identity, probabilistic weight, and inferential origin
are intentionally outside this first gate.

\subsection{Gate 2: the topological exchange boundary}
\label{sec:matroid-topological-boundary}

The repository also contains a topological evidence semantics in which an
algebraic interior operator induces a closure \(\operatorname{cl}\). That closure
is extensive, monotone, idempotent, and preserves binary unions. The last fact
makes the interaction with matroid exchange unusually rigid. Let
\[
  x\preceq_{\operatorname{cl}} y
  \quad\Longleftrightarrow\quad
  x\in\operatorname{cl}(\{y\}).
\]
For a genuine point-set topology this is the specialization preorder (up to the
usual orientation convention). Symmetry of this preorder is the standard
\(R_0\) separation condition \cite{Willard1970}.

\begin{theorem}[Topological exchange characterization]
For the topological closure generated by an \texttt{InteriorSemantics}, the
following are equivalent:
\begin{enumerate}[label=(\roman*)]
  \item the closure satisfies the Mac Lane--Steinitz exchange condition;
  \item singleton specialization is symmetric:
  \[
    x\in\operatorname{cl}(\{y\})
    \quad\Longrightarrow\quad
    y\in\operatorname{cl}(\{x\}).
  \]
\end{enumerate}
Thus, for ordinary topological spaces, the repository-level matroid exchange
boundary is exactly the \(R_0\) boundary.
\textup{\textsc{Status: Lean-proved in
\texttt{PEL4/MatroidTopologicalBridge.lean}; declaration-level audit included.}}
\end{theorem}
\begin{proof}[Proof sketch]
Since topological closure preserves binary unions,
\[
  \operatorname{cl}(A\cup\{y\})
  =\operatorname{cl}(A)\cup\operatorname{cl}(\{y\}).
\]
If \(x\) enters the closure only after adjoining \(y\), then
\(x\in\operatorname{cl}(\{y\})\); singleton symmetry gives
\(y\in\operatorname{cl}(\{x\})\), hence
\(y\in\operatorname{cl}(A\cup\{x\})\). Conversely, exchange over the empty
base sends \(x\in\operatorname{cl}(\{y\})\) to
\(y\in\operatorname{cl}(\{x\})\), because the topological closure of the empty
set is empty.
\end{proof}

This result distinguishes two forms of ``being generated by nearby evidence.''
Topological closure may encode directional approximation: one point can lie in
the closure of another without reciprocal dependence. Matroid exchange forbids
that asymmetry at the singleton level. In a genuine \(T_0\) topology, adding the
exchange condition therefore pushes the singleton specialization relation to
equality, i.e. to the \(T_1\) boundary.

\begin{proposition}[Sierpiński obstruction]
The two-point Sierpiński topology already used as a finite witness in the 4-PEL
topological semantics does not satisfy matroid exchange. At its distinguished
points,
\[
  \mathsf{focus}\in\operatorname{cl}(\{\mathsf{neighbour}\}),
  \qquad
  \mathsf{neighbour}\notin\operatorname{cl}(\{\mathsf{focus}\}).
\]
Hence singleton specialization is asymmetric, so the topological closure fails
exchange.
\textup{\textsc{Status: Lean-proved in
\texttt{PEL4/MatroidTopologicalSierpinski.lean}.}}
\end{proposition}

The obstruction is internal to the existing 4-PEL topological research line;
it is not an unrelated topology introduced solely to refute exchange. The same
Sierpiński witness previously exposed directional boundary behaviour for
\(\Box\) and \(\Diamond\), and Gate 2 now identifies that directionality as the
precise reason the closure is not matroidal.

\begin{proposition}[Channel topology realization]
The coarse Gate-1 positive/negative channel matroid is itself topological.
Define
\[
  \operatorname{Int}_{\chi}(A)
  =\{e : \forall a\,(\chi(a)=\chi(e)\Rightarrow a\in A)\}.
\]
This is the partition topology generated by the polarity classes, and its dual
closure satisfies
\[
  \operatorname{cl}_{\operatorname{Int}_{\chi}}(A)
  =\operatorname{cl}_{\chi}(A).
\]
Its singleton specialization relation is symmetric and therefore satisfies
exchange.
\textup{\textsc{Status: Lean-proved in
\texttt{PEL4/MatroidTopologicalBridge.lean}.}}
\end{proposition}

Gate 2 therefore locates the original channel matroid inside the intersection
of the two structural programs: it is simultaneously a partition-topological
closure and a matroidal closure. The Sierpiński semantics lies outside that
intersection. This gives the 4-PEL evidence layer a first nontrivial structural
boundary rather than merely two parallel mathematical descriptions.

\paragraph{Formal scope.}
The repository's \texttt{MatroidClosure} is intentionally a lightweight
closure-axiom presentation and currently does not add a finite-character axiom.
The characterization above is therefore a theorem about exactly that formal
notion. On the finite carriers used by the explicit 4-PEL witnesses, finite
character is automatic; infinite-carrier finitarity remains a separate
extension.

\paragraph{Next structural questions.}
The next natural gate is combinatorial rather than topological: formalize
matroid circuits as minimal dependent evidence sets and determine what a
minimal evidential redundancy means in 4-PEL. After that, deletion and
contraction can be tested as formal models of removing an evidence source and
quotienting by accepted background evidence. A further question is whether the
six-cell reliability refinement preserves the \(R_0\)-like exchange boundary
under coarse projection.
